% ---------------------------------------------------------------------------
% Response to R2 Major 3 -- JEDC revision
% Body only: paste after the previous comment in the combined response file.
% Needs booktabs. Top heading is \section (numbered), subheadings \subsection*.
% ---------------------------------------------------------------------------

\section{R2, Major 3}
{\itshape Oil imports during default}

\begin{quote}\itshape
In the model, a country in default keeps importing oil at the world price with
no friction. The only direct cost of default is the endowment loss $h(y)$.
[\ldots] The author should discuss this assumption; even better would be a
robustness exercise in which default also disrupts oil imports, for example
through a price wedge or partial exclusion from oil purchases during autarky.
\end{quote}

We have done both. The assumption is discussed below and the robustness
exercise has been run, since it is a parameter of the model already. The result
is negative, and we report it in full because the way it fails is informative
about the mechanism.

\subsection*{The assumption}

Oil is bought on spot markets, usually by private importers and refiners rather
than by the sovereign, and a sovereign default does not by itself bar a
country's firms from purchasing crude. The channel through which a default
episode disrupts fuel supply is the loss of foreign exchange and of trade
credit --- the inability to open letters of credit --- rather than exclusion
from the oil market as such. The paper's motivating example works that way: the
2022 fuel crisis was a balance-of-payments and credit event.

That makes the frictionless assumption a defensible simplification rather than a
neutral one, and the referee is right that the paper should say so. Mendoza and
Yue (2012) make the loss of imported inputs the central cost of default, so a
reader coming from that literature will expect the channel to be present and its
absence should be a stated modelling choice.

\subsection*{The exercise}

While in default or excluded, the country pays $p(1+\kappa)$ for imported oil.
This is a parameter of the model, and $\kappa = 0$ reproduces the baseline
exactly.

\begin{table}[h]
\centering\footnotesize
\setlength{\tabcolsep}{5pt}
\begin{tabular}{cccccc}
\toprule
$\kappa$ & mean spread & debt/GDP & default \% p.a. & \multicolumn{2}{c}{spread} \\
\cmidrule(lr){5-6}
         &             &          &                 & low vol.\ & high vol.\ \\
\midrule
0    & 3.59\% & 26.1\% & 3.35\% & 3.27\% & 4.05\% \\
0.10 & 2.94\% & 29.7\% & 2.80\% & 2.70\% & 3.30\% \\
0.25 & 2.01\% & 36.7\% & 1.81\% & 1.92\% & 2.14\% \\
0.50 & 1.30\% & 48.2\% & 1.27\% & 1.30\% & 1.29\% \\
\bottomrule
\end{tabular}
\end{table}

\noindent
$\sigma_{\text{oil}} = 0.25$, every other parameter at its baseline value.
Spreads are whole-simulation means and debt is relative to annual output. The
levels are therefore not the Table 3 statistics, which average over the hundred
quarters preceding each default episode: on that basis the baseline
economy has a mean spread of 4.08\% and debt/GDP of 34.3\%, against targets of
4.1\% and 34\%. Every column is computed the same way at every $\kappa$, so the
comparison across rows is like for like.

Three things move the wrong way, all monotonically.

\textbf{The state dependence collapses.} The gap between high- and
low-volatility spreads falls from $+0.77$ to $+0.22$ and then to zero: at
$\kappa = 0.50$ the simulated spread is the same in the most and least volatile
states to two decimal places. The paper's central prediction is not weakened at
the margin, it is gone.

\textbf{The spread level falls by nearly two thirds}, 3.59\% to 1.30\%. A
more painful default means less default is priced, and the baseline's match to
the 4.1\% spread target is lost.

\textbf{Debt runs away}, rising by more than four fifths across the range, for
the same reason: a costlier default sustains more borrowing. The baseline
already matches the paper's 34\% debt target, so every positive $\kappa$ pushes
the debt moment above it.

These are not independent, and that is the point. The baseline is the only
version of the model that matches the debt target, and it is also the version
with the full state dependence. A disruption premium moves both away at once:
debt rises above target while the volatility gap shrinks, and there is no
intermediate value at which the calibration is recovered. \textbf{There is no
disruption premium that buys the mechanism without costing the debt moment.}

\subsection*{Why it fails, which is the useful part}

The extra cost of defaulting is $\kappa pm$: it is proportional to the oil
price. A dearer barrel therefore does two things at once. It squeezes the
importer, which pushes it toward default --- the paper's mechanism. And it
raises the penalty for defaulting, which pushes it away. The two offset, and
past roughly $\kappa = 0.2$ the penalty dominates.

This is not a numerical artefact. It is the cancellation that makes the oil
price a weak driver of default incentives whenever the shock is near-permanent:
a persistent price rise shifts the value of repaying and the value of defaulting
almost equally. A disruption premium does not break that symmetry, it tilts it
the wrong way, because the premium is largest exactly when the price is highest.

\subsection*{What would work, and why we have not done it}

For a price rise to raise the spread more strongly, the default penalty must
\emph{fall} when oil is dear, so that defaulting becomes relatively more
attractive precisely when the importer is squeezed. The natural candidate is to
express the penalty as a share of value added --- output net of the oil bill ---
rather than of the endowment, since value added shrinks when oil is expensive.

We have not implemented it. The current penalty is in endowment units, and
tying it to value added introduces a fixed point, since the penalty affects
productive capacity, which affects oil demand, which affects value added. That
requires an iterative solve and is a structural change to the production and
default block rather than a robustness check.

\subsection*{Suggested manuscript wording}

\begin{quote}
The model assumes that a country in default continues to import oil at the world
price. The channel through which default episodes disrupt fuel supply in
practice is the loss of foreign exchange and trade credit rather than exclusion
from oil markets, so we treat this as a simplification and report what relaxing
it does. Charging a defaulting country $p(1+\kappa)$ for oil raises the cost of
default in proportion to the oil price, which both squeezes the importer and
raises the penalty for defaulting; past $\kappa \approx 0.2$ the second effect
dominates. The gap between high- and low-volatility spreads falls from 0.77
percentage points to zero as $\kappa$ rises to 0.50, while debt-to-GDP rises by
more than four fifths above a baseline that already matches its 34 per cent
target, so the disruption premium cannot be calibrated to match the debt moment
and preserve the state dependence at the same time. A penalty proportional to
value added rather than to the endowment would work in the required direction,
but it introduces a fixed point between the default cost and oil demand, and we
leave it to future work.
\end{quote}
